% whitepaper.tex -- the Passiflora polDisc pre-filter: what it computes, why it is correct, how that was checked, what it still assumes.
% Self-contained: the proofs are in this document.  Compile:  pdflatex --disable-installer -interaction=nonstopmode -halt-on-error whitepaper.tex  (twice)
\documentclass[10pt,a4paper]{article}
\usepackage[T1]{fontenc}
\usepackage{lmodern}
\usepackage[margin=2.5cm]{geometry}
\usepackage{amsmath,amssymb,amsthm}
\usepackage{booktabs,tabularx}
\usepackage{enumitem}
\usepackage{xcolor}
\usepackage{xspace}
\usepackage[colorlinks=true,linkcolor=blue!50!black,citecolor=blue!50!black,urlcolor=blue!50!black]{hyperref}

\newcommand{\F}{\mathbb{F}}
\newcommand{\Fq}{\mathbb{F}_q}
\newcommand{\Z}{\mathbb{Z}}
\newcommand{\Q}{\mathbb{Q}}
\newcommand{\K}{\mathbb{K}}
\DeclareMathOperator{\Res}{Res}
\DeclareMathOperator{\disc}{disc}
\DeclareMathOperator{\lc}{lc}
\DeclareMathOperator{\cont}{cont}
\newcommand{\leg}[2]{\left(\frac{#1}{#2}\right)}
\newcommand{\code}[1]{\texttt{#1}}
\newcommand{\stproved}{\textsf{\small[proved]}\xspace}
\newcommand{\stquoted}{\textsf{\small[standard result, quoted]}\xspace}
\newcommand{\stnum}{\textsf{\small[checked numerically]}\xspace}
\newcommand{\streport}{\textsf{\small[reported, not re-checked]}\xspace}
\newcommand{\stassume}{\textsf{\small[named assumption]}\xspace}

\theoremstyle{plain}
\newtheorem{proposition}{Proposition}
\newtheorem{lemma}[proposition]{Lemma}
\newtheorem{theorem}[proposition]{Theorem}
\newtheorem{corollary}[proposition]{Corollary}
\newtheorem*{thmA}{Theorem A}
\newtheorem*{thmB}{Theorem B}
\newtheorem*{thmC}{Theorem C}
\newtheorem*{thmD}{Theorem D}
\newtheorem*{thmE}{Theorem E}
\theoremstyle{definition}
\newtheorem{definition}[proposition]{Definition}
\newtheorem{remark}[proposition]{Remark}

\setlist{itemsep=2pt,topsep=4pt}
\setlength{\parskip}{3pt}
\setlength{\emergencystretch}{3em}
\renewcommand{\arraystretch}{1.12}

\begin{document}
\begin{center}
{\LARGE An exact integer pre-filter for the NumberFields@home discriminant test}\\[8pt]
{\large The Passiflora polDisc pre-filter: what it computes, why it is correct, and how that was checked}\\[6pt]
2 November 2026
\end{center}

\begin{abstract}
NumberFields@home searches for number fields of degree $10$ by testing a very large number of monic integer polynomials with an exact discriminant test, which
dominates the running time of its CPU application, GetDecics.  Passiflora places a pre-filter in front of that test.  The pre-filter discards a polynomial
when the Legendre symbol of its discriminant is $-1$ modulo one of $30$ primes chosen for the work unit; such a polynomial provably cannot pass.  A pre-filter
of this kind is acceptable only if it never discards a polynomial that passes.  This paper describes an implementation in integer arithmetic (signed Montgomery
reduction on the CPU, unsigned Montgomery reduction on the GPU) and proves that the program as written makes exactly the decisions of the exact algorithms,
independently of the floating-point environment.  It states the contract between the filter and the application and reports the evidence: an exhaustive test of
the compiled filter on $25.5$ billion polynomials at tiny primes, an exact oracle with deliberate faults, a differential test of GetDecics' own exact test, the
$31$ test work units in every build, and the GPU kernel on the CPU and on a device.  It also proves what GetDecics' exact test computes, under one condition on the
data that the search regions are designed to meet: the test's first $64$-bit steps do not overflow.  Every result carries its status: proved, checked
numerically, quoted, or assumed.
\end{abstract}

\tableofcontents

%======================================================================
\section{Introduction}

NumberFields@home searches for number fields of degree $10$ by enumerating monic polynomials $f=x^{10}+a_9x^9+\dots+a_0$ with integer coefficients.  The
program that volunteers run is \emph{GetDecics} \cite{gd}.  Its final step is an exact \emph{discriminant test}, \code{polDiscTest\_cpu}, built on GMP:

\begin{definition}[Passing]\label{def:pass}
Let $S$ be a finite set of primes (read from the work unit; $S=\{2,5\}$ in the current searches and in every test work unit).  A monic $f\in\Z[x]$ \emph{passes} if
$D=\disc(f)\neq0$ and $|D|=\bigl(\prod_{p\in S}p^{\alpha_p}\bigr)m^2$ for some integers $\alpha_p\ge0$, $m\ge1$.
\end{definition}

Almost no polynomial passes, and the stock CPU application spends about $98\%$ of its time in this test \streport.  \emph{Passiflora} is a patch to GetDecics
(source at commit \code{638ca4e}) with new source files for a pre-filter and for the GPU code.  The enumeration and the exact test remain GetDecics' own code;
Passiflora decides which polynomials reach the test, and its output must be exactly that of the stock program.

\paragraph{The idea.}  If $f$ passes, then modulo a suitable prime $q$ its discriminant is a square or zero (Proposition~\ref{prop:rule}).  So a Legendre
symbol $\leg{\disc f}{q}=-1$ at one such prime proves that $f$ fails, and the filter computes these symbols.  The CPU filter uses $30$ primes: $24$ \emph{small}
ones below $2^{13}$ ($8161,8089,\dots,6481$ for $S=\{2,5\}$) and six \emph{large} ones below $2^{25}$ ($33554249,\dots,33553649$).  Each rejects about half of
the polynomials that reach it; the survivors of all $30$ go to the exact test, which decides them as before.  The GPU kernel applies the same filter at the
$24$ smallest suitable primes in $(256,8192)$ ($281,\dots,1489$ for $S=\{2,5\}$, none of them the CPU's) before it hands its survivors to the CPU.

\paragraph{Why exactness needs care.}  The filter must never reject a polynomial that passes, or the project loses fields.  An earlier version stored
residues in floating point and reduced with the rounding trick $x-q\cdot((x\,q^{-1}+M)-M)$, which is exact only if every operation is rounded to nearest in
the declared precision.  On the real application that assumption fails in practice \stnum: built for x87 arithmetic (as a 32-bit x86 build is), the application
lost valid polynomials in four of the seven smallest test work units ($6$, $12$, $6$ and $5$), among them the only real result of one unit; a library that switches
the rounding mode to ``up'' before the filter is initialised makes it reject valid polynomials.  A filter whose correctness depends on the floating-point
environment of a volunteer's machine cannot be proved correct for the program that ships.  The filter described here uses integer arithmetic only.

\paragraph{Results.}  Section~\ref{sec:background} proves the rejection rule.  Section~\ref{sec:chains} gives the two Euclidean chains that compute the symbol
and proves them exact over $\Fq$ (Theorems A, B, C).  Section~\ref{sec:integer} proves that the integer program computes exactly those decisions, for every input
(Theorem D), and derives the size limits of its primes.  Section~\ref{sec:gpu} does the same for the GPU kernel (Theorem E).  Section~\ref{sec:poldisc} proves
what GetDecics' exact test computes.  Section~\ref{sec:contract} states the contract with the application, Section~\ref{sec:evidence} the evidence, and
Section~\ref{sec:speed} the performance.  Section~\ref{sec:ledger} lists what everything rests on.  Each result is marked \stproved, \stnum (by a complete or
stated finite computation, or by tests), \stquoted, \streport or \stassume.

\paragraph{One assumption about the data.}  Everything is proved for every input except one link: that GetDecics' exact test computes the definition of
passing.  It does so only while its first two iterations, which it runs in $64$-bit integers without overflow checks, stay in range: hypothesis (H) of
Section~\ref{sec:poldisc}.  The search regions keep these values below $2^{60}$, and none of the $14.8$ billion polynomials of the $31$ test work units
overflows.  For a polynomial outside, GetDecics can give a wrong verdict, and Passiflora's counters can then differ from the stock program's.  Whoever widens a
search region, or gives the program other polynomials, has to check (H) first; the source says so where the filter is called and above the $64$-bit steps.

%======================================================================
\section{The rejection rule}\label{sec:background}

Let $\F$ be a field and fix an algebraically closed field $\K\supseteq\F$.  For nonzero $V,P\in\F[x]$ with $V=\lc(V)\prod_{i=1}^{v}(x-\beta_i)$ over $\K$, the
\emph{resultant} is $\Res(V,P)=\lc(V)^{\deg P}\prod_{i}P(\beta_i)$, and $\Res(V,0)=0$ for $\deg V\ge1$.  It lies in $\F$ \stquoted \cite[Chapter~IV]{lang}.  Directly
from the definition \stproved:
\begin{enumerate}[label=(R\arabic*)]
\item $\Res(c,P)=c^{\deg P}$ and $\Res(V,c)=c^{\deg V}$ for a nonzero constant $c$;
\item $\Res(U,V)=(-1)^{\deg U\deg V}\Res(V,U)$, because $\prod_{i,j}(\alpha_j-\beta_i)$ and $\prod_{i,j}(\beta_i-\alpha_j)$ differ in $\deg U\deg V$ signs;
\item if $\deg V\ge1$ and $R=P\bmod V$, then $\Res(V,P)=0$ when $R=0$, and $\Res(V,P)=\lc(V)^{\deg P-\deg R}\Res(V,R)$ otherwise, because $P(\beta_i)=R(\beta_i)$.
\end{enumerate}
The \emph{discriminant} of a monic $f$ of degree $n$ with roots $\rho_1,\dots,\rho_n$ is $\disc f=\prod_{i<j}(\rho_i-\rho_j)^2$.

\begin{proposition}\label{prop:F1}\stproved
Let $f\in\F[x]$ be monic of degree $n\ge1$ with $n\cdot1\neq0$ in $\F$.  Then $\deg f'=n-1$, $\lc(f')=n$, and $\disc f=(-1)^{n(n-1)/2}\Res(f,f')=(-1)^{n(n-1)/2}\Res(f',f)$.
\end{proposition}
\begin{proof}
$f'(\rho_i)=\prod_{j\neq i}(\rho_i-\rho_j)$ by the product rule, so $\Res(f,f')=\prod_if'(\rho_i)=\prod_{i<j}(\rho_i-\rho_j)(\rho_j-\rho_i)=(-1)^{n(n-1)/2}\disc f$.
By (R2), $\Res(f',f)=(-1)^{n(n-1)}\Res(f,f')=\Res(f,f')$.
\end{proof}

\begin{proposition}\label{prop:reduction}\stproved
Let $f\in\Z[x]$ be monic of degree $n\ge2$ and $q$ a prime with $q\nmid n$.  Then $\disc f\in\Z$, and its residue modulo $q$ is the discriminant of $f\bmod q$.
\end{proposition}
\begin{proof}
$\Res(f,f')$ is the determinant of the Sylvester matrix \stquoted \cite[Section~3.3]{cohen}, a polynomial with integer coefficients in the coefficients of $f$
and $f'$.  As $q\nmid n$, $f'\bmod q$ still has degree $n-1$, so the Sylvester matrix over $\Fq$ is the reduction of the one over $\Z$, and the determinant
commutes with reduction.  Apply Proposition~\ref{prop:F1} on both sides.
\end{proof}

For an odd prime $q$ write $\chi(a)=\leg aq$.  \emph{Euler's criterion} $a^{(q-1)/2}\equiv\chi(a)\pmod q$ \stquoted gives: $\chi$ is multiplicative;
$\chi(c^2a)=\chi(a)$ for $c\not\equiv0$; $\chi(c^e)=\chi(c)^{e\bmod2}$ for every integer $e$ (negative $e$ meaning powers of $c^{-1}$), so only the parity of an
exponent matters; and $\chi(-1)=(-1)^{(q-1)/2}$, so $-1$ is a square exactly when $q\equiv1\pmod4$.

\begin{proposition}[The rejection rule]\label{prop:rule}\stproved
Let $f$ pass for $S$, and let $q$ be an odd prime with $q\equiv1\pmod4$ at which every $p\in S$ is a nonzero square.  Then $\leg{\disc f}{q}\in\{0,1\}$.  So if
$\leg{\disc f}{q}=-1$ at such a prime, $f$ does not pass.
\end{proposition}
\begin{proof}
$D=\pm\bigl(\prod_pp^{\alpha_p}\bigr)m^2$, so $\chi(D)=\chi(\pm1)\prod_p\chi(p)^{\alpha_p}\chi(m)^2=\chi(\pm1)\chi(m)^2\in\{0,1\}$, as $\chi(-1)=1$ and $\chi(p)=1$.
\end{proof}

A prime $q$ with these properties is called \emph{usable} for $S$.  The conditions on $S$ enter only here; the chains below work for any odd prime $q\nmid n$.

%======================================================================
\section{Two Euclidean chains for one Legendre symbol}\label{sec:chains}

Throughout this section $q$ is an odd prime, $f\in\Fq[x]$ is monic of degree $n\ge2$ with $q\nmid n$, and arithmetic is exact in $\Fq$.  The filter computes the
symbol of $\Res(f',f)$ by Euclid's algorithm on $f$ and $f'$, without inverses: instead of dividing by a leading coefficient it multiplies by it.

\begin{lemma}[Pseudo-division step]\label{lem:L1}\stproved
Let $U,V\in\Fq[x]$ be nonzero, $u=\deg U\ge v=\deg V$, $l_u=\lc U$, $l_v=\lc V$, and $W=l_vU-l_ux^{u-v}V$, so $\deg W<u$.  Then $W(\beta)=l_vU(\beta)$ at every
root $\beta$ of $V$; if $W\neq0$, then $\Res(V,W)=l_v^{\,e}\Res(V,U)$ with $e=\deg W-u+v$; and if $W=0$ and $v\ge1$, then $\Res(V,U)=0$.
\end{lemma}
\begin{proof}
$V(\beta)=0$ gives the first claim.  Then $\Res(V,W)=l_v^{\deg W}\prod_iW(\beta_i)=l_v^{\deg W+v}\prod_iU(\beta_i)=l_v^{\deg W+v-u}\Res(V,U)$.  If $W=0$, then
$U(\beta_1)=0$.
\end{proof}

By the remarks on $\chi$, a step with $e$ even does not change the symbol, and a step with $e$ odd changes it by $\chi(l_v)$.

\begin{definition}
The \emph{remainder sequence} of $(f,f')$ is $r_0=f$, $r_1=f'$, $r_{s+1}=r_{s-1}\bmod r_s$ while $r_s\neq0$ and $\deg r_s\ge1$.  It is \emph{normal} if
$\deg r_s=n-s$ for $s=0,\dots,n-1$ (then $r_n$ is a constant, possibly $0$), and \emph{abnormal} otherwise.
\end{definition}

\subsection{\code{chain\_block}, the fast chain}

\code{chain\_block} stores coefficients by increasing powers ($P[k]$ is the coefficient of $x^k$) and runs $64$ polynomials in the same instructions.  For
$s=1,\dots,n-1$, with $m=n-s$, $A=P_{s-1}$, $B=P_s$ ($P_0=f$, $P_1=f'$): $\lambda_s=B[m]$, $\tau=A[m+1]$; round 1, $R=\lambda_sA-\tau xB$; round 2,
$P_{s+1}=\lambda_sR-R[m]\,B$; and, if $m\ge2$, the check $\code{ok}\leftarrow\code{ok}\wedge(P_{s+1}[m-1]\neq0)$.  The coefficients that are not computed are zero:
the coefficient of $x^{m+1}$ in $R$ is $\lambda_s\tau-\tau\lambda_s$, and that of $x^m$ in $P_{s+1}$ is $\lambda_sR[m]-R[m]\lambda_s$.  At the end
$\code{base}=P_n[0]$, $\code{acc}=\code{base}^{(q-1)/2}$, and the polynomial is rejected when \code{ok} holds and $\code{acc}=-1$.

\begin{lemma}\label{lem:L2}\stproved
Let $B$ have degree exactly $m\ge1$, $\lambda=\lc B$, $\deg A\le m+1$, $\tau$ the coefficient of $x^{m+1}$ in $A$, $R=\lambda A-\tau xB$, $t=R[m]$ and
$R_2=\lambda R-tB$.  Then $R_2=\lambda^2(A\bmod B)$.
\end{lemma}
\begin{proof}
$\lambda^2A=\lambda R+\lambda\tau xB=(\lambda\tau x+t)B+R_2$ with $\deg R_2<m$; by uniqueness of division $R_2=(\lambda^2A)\bmod B=\lambda^2(A\bmod B)$.
\end{proof}

\begin{lemma}\label{lem:sigma}\stproved
Put $\sigma_0=\sigma_1=1$ and $\sigma_{s+1}=\lambda_s^2\sigma_{s-1}$.  If $\lambda_1,\dots,\lambda_S\neq0$ ($1\le S\le n-1$), then $\deg r_s=n-s$ for $1\le s\le S$,
$P_s=\sigma_sr_s$ for $0\le s\le S+1$, and every $\sigma_s$ is a nonzero square.
\end{lemma}
\begin{proof}
Induction on $s$.  If $P_{s-1}=\sigma_{s-1}r_{s-1}$ and $P_s=\sigma_sr_s$ with $\sigma$'s nonzero and $\lambda_s\neq0$, then $\deg P_s=n-s$, so $\deg r_s=n-s$, and
Lemma~\ref{lem:L2} gives $P_{s+1}=\lambda_s^2(\sigma_{s-1}r_{s-1}\bmod\sigma_sr_s)=\lambda_s^2\sigma_{s-1}(r_{s-1}\bmod r_s)=\sigma_{s+1}r_{s+1}$.
\end{proof}

\begin{proposition}[The check detects exactly abnormality]\label{prop:okiffnormal}\stproved
At the end of \code{chain\_block}, \code{ok} is true if and only if the remainder sequence of $(f,f')$ is normal; if the sequence first becomes abnormal at
$r_{s+1}$, the check first fails at step $s$.
\end{proposition}
\begin{proof}
$\lambda_1=n\neq0$, and the check of step $s$ tests $\lambda_{s+1}$, so \code{ok} holds iff $\lambda_2,\dots,\lambda_{n-1}\neq0$; then Lemma~\ref{lem:sigma} gives
$\deg r_s=n-s$ for all $s\le n-1$.  Conversely, if the sequence is normal and $\lambda_s=0$ for a smallest $s\ge2$, Lemma~\ref{lem:sigma} with $S=s-1$ gives
$\lambda_s=\sigma_s\lc(r_s)\neq0$, a contradiction.  The same argument at the first abnormal index gives the last claim.
\end{proof}

\begin{lemma}\label{lem:resnormal}\stproved
If the remainder sequence is normal, $\Res(f',f)=\bigl(\prod_{s=1}^{n-1}\lc r_s\bigr)^2r_n$.
\end{lemma}
\begin{proof}
For $s\le n-2$, (R3) and (R2) give $\Res(r_s,r_{s-1})=\lc(r_s)^2\Res(r_s,r_{s+1})=\lc(r_s)^2\Res(r_{s+1},r_s)$, as $(n-s)(n-s-1)$ is even: a normal sequence
produces no signs.  For $s=n-1$, $\Res(r_{n-1},r_{n-2})=\lc(r_{n-1})^2r_n$ by (R3) and (R1) (both sides are $0$ if $r_n=0$).  Multiply.
\end{proof}

\begin{thmA}[\code{chain\_block}]\stproved
If \code{ok} is true at the end, then $\code{base}=\sigma_nr_n$ and $\chi(\code{base})=\chi\bigl((-1)^{n(n-1)/2}\bigr)\chi(\disc f)$.  For $q\equiv1\pmod4$
this is $\chi(\disc f)$, $\code{acc}=\chi(\disc f)$, and a rejection means $\chi(\disc f)=-1$.
\end{thmA}
\begin{proof}
By Proposition~\ref{prop:okiffnormal} and Lemma~\ref{lem:sigma}, $\code{base}=P_n[0]=\sigma_nr_n$ with $\sigma_n$ a nonzero square, so $\chi(\code{base})=\chi(r_n)$.  By
Proposition~\ref{prop:F1} and Lemma~\ref{lem:resnormal}, $\disc f=(-1)^{n(n-1)/2}\Pi^2r_n$ with $\Pi\neq0$.  For $q\equiv1\pmod4$, $\chi(-1)=1$; Euler's criterion gives
$\code{acc}=\chi(\code{base})$.
\end{proof}

For $n=10$ the sign $(-1)^{45}$ is $-1$, so the condition $q\equiv1\pmod4$ is needed: with $q\equiv3\pmod4$, \code{chain\_block} returns the wrong symbol for every
polynomial with a normal sequence and $\disc f\not\equiv0$.

\subsection{\code{chain\_general}, the general chain}

\code{chain\_general} stores polynomials \emph{left-aligned}: entry $i$ of $P$ is the coefficient of $x^{\deg P-i}$, and the entries beyond the degree are zero.
Then one elementary step $W=l_vU-l_ux^{u-v}V$ is the same formula in every lane, whatever the degrees: $W(i)=l_vU(i+1)-l_uV(i+1)$ for $0\le i<u$ is $W$ left-aligned
with nominal degree $u-1$ (the coefficient of $x^{u-j}$ in $x^{u-v}V$ is $V(j)$).  One lane runs:
\begin{enumerate}[label=(\alph*)]
\item $U=f$, $V=f'$, $\code{mm}=1$; in each iteration form $W$ and shift out its leading zeros, $w=\deg W$;
\item if $W=0$, stop: the polynomial is kept;
\item with $e=w-u+v$, if $e$ is odd, $\code{mm}\leftarrow\code{mm}\cdot l_v$;
\item if $w\ge v$, $U\leftarrow W$; otherwise $(U,V)\leftarrow(V,W)$ (a \emph{swap});
\item if now $V=c_0$ is a constant: if $\deg U$ is odd, $\code{mm}\leftarrow\code{mm}\cdot c_0$; stop.
\end{enumerate}
Then $\code{acc}=\code{mm}^{(q-1)/2}$, and the polynomial is rejected when the lane stopped at (e) and $\code{acc}=-1$.

\begin{proposition}[Invariant and termination]\label{prop:invariant}\stproved
At the start of every iteration $U,V\neq0$, $u\ge v\ge1$, and there are $\varepsilon\in\{\pm1\}$ and $\gamma\in\Fq^\times$ with
\[ \Res(f',f)=\varepsilon\cdot\code{mm}\cdot\gamma^2\cdot\Res(V,U), \tag{I}\]
where $\varepsilon=(-1)^N$, $N$ the sum of $vw$ over the swaps so far.  The lane stops within $2n-2$ iterations.
\end{proposition}
\begin{proof}
At the start $U=f$, $V=f'$ ($l_v=n\neq0$), $\code{mm}=1$.  In an iteration that goes on, $W\neq0$ and Lemma~\ref{lem:L1} gives $\Res(V,U)=l_v^{-e}\Res(V,W)$; write
$l_v^{-e}=\delta^2$ ($e$ even) or $l_v\delta^2$ ($e$ odd).  Step (c) multiplies \code{mm} by $l_v$ exactly in the second case, so (I) holds for $(W,V)$ with
$\gamma\delta$; after a swap (R2) adds the sign $(-1)^{vw}$.  The new pair satisfies $u\ge v$, and the loop goes on only with $v\ge1$.  Each iteration that goes on
lowers $u+v$ by at least one; it starts at $2n-1$ and is at least $2$, so there are at most $2n-2$ iterations.
\end{proof}

\begin{thmB}[\code{chain\_general}]\stproved
The lane stops within $2n-2$ iterations ($18$ for $n=10$; the code allows $40$).  If it stops at (b), then $\disc f=0$ and the polynomial is kept.  If it stops at
(e), then $\code{mm}\neq0$ and $\chi(\code{mm})=\chi\bigl((-1)^{n(n-1)/2}\varepsilon\bigr)\chi(\disc f)$.  For $q\equiv1\pmod4$ this is $\chi(\disc f)$,
$\code{acc}=\chi(\disc f)$, and a rejection means $\chi(\disc f)=-1$.
\end{thmB}
\begin{proof}
At (b), (I) holds and $\Res(V,U)=0$ by Lemma~\ref{lem:L1}, so $\Res(f',f)=0$.  At (e), (I) holds for the new pair with $V=c_0\neq0$, and $\Res(V,U)=c_0^{\deg U}$ by
(R1), whose odd part step (e) puts into \code{mm}; so $\Res(f',f)=\varepsilon\cdot\code{mm}\cdot\gamma'^2$ with \code{mm} a product of nonzero leading coefficients.
Apply Proposition~\ref{prop:F1} and $\chi(\pm1)=1$ for $q\equiv1\pmod4$.
\end{proof}

The swap signs are why $q\equiv1\pmod4$ matters here too \stnum: with primes $q\equiv3\pmod4$, \code{chain\_general} was wrong in about a quarter of the degree-$4$
cases and $70\%$ of the degree-$10$ cases tested.

\subsection{The filter}\label{sec:filter}

\code{pdf\_init(numP, pSet)} chooses the primes: a prime is \emph{chosen} if it is usable for $S$ and lies in the proved range (the \emph{guard},
Section~\ref{sec:limits}): the $24$ largest below $2^{13}$ and the six largest below $2^{25}$.  \code{pdf\_filter} then runs every surviving polynomial, stage by
stage, through \code{chain\_block} (pass 1), re-runs the lanes whose check failed through \code{chain\_general} (pass 2), and drops a polynomial as soon as a chain
rejects it.  A stage at which more than half of the lanes left \code{chain\_block} in the last call sends every lane to \code{chain\_general} directly
(\code{prefgen}), and every $64$th call probes \code{chain\_block} again.  Survivors of all stages get flag $1$, all others $0$; a polynomial whose leading
coefficient is not $1$ always gets flag $1$.

\begin{thmC}[soundness and exactness]\stproved
In exact arithmetic, for every input, block size and routing state: \code{pdf\_filter} gives $f$ the flag $0$ if and only if $f$ is monic and
$\leg{\disc f}{q_j}=-1$ at some stage $j$, and then $f$ is dropped at the first such stage.  In particular it never rejects a polynomial that passes.
\end{thmC}
\begin{proof}
A flag $0$ needs a leading coefficient $1$, and the chains read $a_0,\dots,a_9$ with $a_{10}=1$, so they see $f$.  At a stage it reaches, $f$ takes one of three routes:
\code{chain\_general} in pass 1 (\code{prefgen}); \code{chain\_block} with \code{ok} false, then \code{chain\_general} in pass 2; or \code{chain\_block} with \code{ok}
true.  The prime is usable and at least $11$, so $q\nmid10$.  By Theorems A and B, on every route $f$ is dropped exactly when $\chi_j(\disc f)=-1$ (on the second route
a kept lane with $\disc f\equiv0$ has symbol $0$).  By Proposition~\ref{prop:reduction} the symbol is that of the integer discriminant, and by
Proposition~\ref{prop:rule} a polynomial with symbol $-1$ does not pass.  Induction over the stages gives the first such stage.
\end{proof}

The converse direction, that the filter rejects \emph{every} polynomial with a symbol $-1$, is what makes a test with an exact oracle possible
(Section~\ref{sec:evidence}).

%======================================================================
\section{Exact integer arithmetic}\label{sec:integer}

\paragraph{Representation.}  Elements of $\Fq$ are integers in a small range around $0$.  For $R=2^k$ let $s_k(x)$ be the residue of $x$ modulo $R$ in
$[-R/2,R/2)$.  The small primes use $R=2^{16}$ in $16$-bit lanes (\code{mont}, $16$ lanes per AVX2 register, compiled to \code{vpmulhw}/\code{vpmullw}); the large
primes use $R=2^{32}$ with $64$-bit products (\code{sredc}).  Each $64$-bit coefficient is split into three $21$-bit pieces and loaded by one reduction
(\code{ldcoef}, \code{ldres}).

\paragraph{What the language must do.}  The proofs use three rules: (S1) an operation whose exact result fits its kind gives that result; (S2) a conversion to a
narrower kind reduces modulo $2^k$; (S3) \code{shifta} is an arithmetic shift and \code{iand} with $2^k-1$ keeps the low $k$ bits.  The Fortran standard leaves
(S2) and the bit model to the processor; the GNU Fortran manual documents (S3) but not (S2), which GCC documents for C.  Standard-conforming replacements exist,
but they made \code{pdf\_filter} $2$ to $6$ times slower in a micro-benchmark ($105$--$126$\,ns per polynomial as written, $572$--$697$ with \code{modulo} and exact
division, $225$--$262$ with offsets; AVX2, an earlier configuration with six small primes) \stnum, because the compiler then no longer uses the $16$-bit
multiply-high instructions.  So the code relies on (S2) and (S3), and they are checked three ways: the compiled reductions against exact models, the whole filter on
complete sets at tiny primes (Section~\ref{sec:evidence}), and at every start of the program (Proposition~\ref{prop:kat}), the only one of the three that runs in
every build.

\subsection{The reductions}

\begin{lemma}[Montgomery reduction]\label{lem:montgen}\stproved
Let $q>1$ be odd, $qq'\equiv1\pmod R$, $x\in\Z$, $m=s_k(xq')$ and $u=(x-mq)/R$.  Then $u$ is an integer, $u\equiv xR^{-1}\pmod q$, and $|u|\le|x|/R+q/2$; in
particular $|x|<qR/2$ implies $|u|<q$.
\end{lemma}
\begin{proof}
$mq\equiv xq'q\equiv x\pmod R$; $uR\equiv x\pmod q$ and $R$ is invertible mod $q$; $|m|\le R/2$.
\end{proof}

The only multiple of $q$ in $(-q,q)$ is $0$, and those in $(-2q,2q)$ are $0$ and $\pm q$: these are the zero tests of the code.

\begin{lemma}[\code{mont}, $R=2^{16}$]\label{lem:mont}\stproved
Let $q$ be odd, $q\le32767$, $q'=s_{16}(q^{-1}\bmod2^{16})$, and $a,b$ $16$-bit values.  \code{mont} computes the $32$-bit product $x=ab$, $\code{lo}=s_{16}(x)$,
$m=s_{16}(\code{lo}\cdot q')$, $h_1=\lfloor x/2^{16}\rfloor$, $h_2=\lfloor mq/2^{16}\rfloor$ and returns $h_1-h_2$.  No operation overflows, the result is the $u$ of
Lemma~\ref{lem:montgen}, and $|u|<q$ whenever $|a|,|b|<2q$ and $q\le2^{13}$.
\end{lemma}
\begin{proof}
$|x|\le2^{30}$ and $|\code{lo}\cdot q'|\le2^{30}$ fit $32$ bits (this needs $q'$ signed); by (S2) $m=s_{16}(xq')$.  By (S3) $h_1,h_2$ are the floors, both at most $2^{14}$ in
absolute value.  As $x\equiv mq\pmod{2^{16}}$, the low halves agree and $h_1-h_2=(x-mq)/2^{16}=u$, with $|u|\le2^{14}+q/2\le32767$.  Finally $(2q-1)^2<2^{15}q$
holds exactly for $q\le2^{13}$.
\end{proof}

\begin{lemma}[\code{sredc}, $R=2^{32}$]\label{lem:sredc}\stproved
Let $q<2^{31}$ be odd, $q'=s_{32}(q^{-1}\bmod2^{32})$, and $x$ a $64$-bit value with $|x|<2^{31}q$.  \code{sredc} returns $u\equiv x2^{-32}\pmod q$ with $|u|<q$, and no
operation overflows.
\end{lemma}
\begin{proof}
$s_{32}(x)q'$ is at most $2^{62}$ in absolute value; $m=s_{32}(xq')$; $|x-mq|<2^{32}q<2^{63}$; $x-mq$ is divisible by $2^{32}$, so the arithmetic shift divides exactly;
$|u|<q$ by Lemma~\ref{lem:montgen}.
\end{proof}

\begin{lemma}[Loading]\label{lem:ldres}\stproved
For every $64$-bit $c$, the pieces $p_0=\code{iand}(c,2^{21}-1)$, $p_1=\code{iand}(\code{shifta}(c,21),2^{21}-1)$, $p_2=\code{shifta}(c,42)$ satisfy
$c=p_0+2^{21}p_1+2^{42}p_2$ with $0\le p_0,p_1<2^{21}$, $-2^{21}\le p_2<2^{21}$.  With $d_i=2^{21i}2^{32}\bmod q$, \code{sredc}$(p_0d_0+p_1d_1+p_2d_2)$ is congruent to $c$
modulo $q$ and lies in $(-q,q)$.
\end{lemma}
\begin{proof}
By (S3) the pieces are the digits of $c$ in base $2^{21}$ with a signed top digit.  The sum has absolute value below $3\cdot2^{21}q<2^{31}q$ and is congruent to $2^{32}c$;
apply Lemma~\ref{lem:sredc}.
\end{proof}

\subsection{Ranges, weights and Euler's criterion}

\begin{proposition}[Ranges]\label{prop:ranges}\stproved
At a stage whose prime $q$ was chosen (so $11\le q\le8385$ for a small and $11\le q\le1{,}097{,}010{,}265$ for a large stage):
(a) at a small stage every input coefficient is in $(-q,q)$, every stored coefficient, leading coefficient and \code{base} in $(-2q,2q)$, \code{mm} and \code{acc} in
$(-q,q)$, every argument of \code{mont} in $(-2q,2q)$, and every difference of two results of \code{mont} in $(-2q,2q)$ without overflow;
(b) at a large stage every value is in $(-q,q)$ and every argument $x$ of \code{sredc} satisfies $|x|\le2(q-1)^2<2^{31}q$, formed without overflow.
So the reductions of Lemmas~\ref{lem:mont} and~\ref{lem:sredc} apply everywhere, and a value is $0$ modulo $q$ exactly when it is $0$ or $\pm q$ (small) or $0$ (large).
\end{proposition}
\begin{proof}
For $q\le2^{13}$ (small) and $q\le2^{30}+1$ (large), by induction over the statements: inputs come from Lemma~\ref{lem:ldres}; the constants $1,\dots,10$ are below
$2q$ as $q\ge11$; \code{mont} of arguments in $(-2q,2q)$ lies in $(-q,q)$ (Lemma~\ref{lem:mont}), so a difference of two results lies in $(-2q,2q)$, below $2^{14}$;
the large-stage arguments are $ka_k$, products $ab$ and differences $ab-cd$ of values in $(-q,q)$, so $|x|\le2(q-1)^2$, which is below $2^{31}q$ for $q\le2^{30}+1$;
all other statements copy, move, merge or zero values.  For $8192<q\le8385$ and $2^{30}+1<q\le1{,}097{,}010{,}265$ the claims follow from
Theorem~\ref{thm:limits}, which bounds every argument by a number below $qR/2$.
\end{proof}

\begin{lemma}[Weights]\label{lem:weights}\stproved
Run a chain exactly over $\Fq$ on $f\bmod q$ and the code on the loaded coefficients.  Both take the same branches, and each polynomial of the code is congruent to
its exact counterpart times $R^{-w}$ for an integer $w\ge0$, the same for all its coefficients; the same holds for the leading coefficients and for \code{mm}.  So
every zero test, degree, swap and \code{ok} agrees with the exact run, and $\chi(\code{base})$ and $\chi(\code{mm})$ are the exact symbols.
\end{lemma}
\begin{proof}
A reduced product of $a$ and $b$ is $\equiv abR^{-1}$, and a reduced difference of products $\equiv(ab-cd)R^{-1}$.  In every formula of the chains the two terms carry the
same weight (in \code{chain\_block}, both terms of $\lambda A[k]-\tau B[k-1]$ have weight $w(A)+w(B)$, and both terms of $\lambda R[k]-tB[k]$ weight $w(R)+w(B)$; in
\code{chain\_general}, both terms of $l_vU(i+1)-l_uV(i+1)$ weight $w(U)+w(V)$), so each result has one weight.  $R^{-w}$ is invertible, so zero tests agree, and the
branches depend only on zero tests.  $R=2^{16}=(2^8)^2$ and $2^{32}=(2^{16})^2$ are squares, so no symbol changes.
\end{proof}

\begin{lemma}[Euler's criterion in the Montgomery domain]\label{lem:eulermont}\stproved
The code starts with $\code{acc}=R\bmod q$ and, for each bit of $(q-1)/2$ from the lowest, multiplies \code{acc} by \code{base} (if the bit is $1$) and squares
\code{base}, all with reduced products.  At the end $\code{acc}\equiv R\chi(b)$, where $b$ is the initial \code{base}; so $\chi(b)=-1$ exactly when \code{acc} is
$\code{negr}=(-R)\bmod q$ or $\code{negr}-q$, which is the code's test.
\end{lemma}
\begin{proof}
With $\beta=bR^{-1}$, the invariants $\code{base}\equiv R\beta^{2^i}$, $\code{acc}\equiv R\beta^{E_i}$ hold, as the reduced product of $R\beta^s$ and $R\beta^t$ is
$R\beta^{s+t}$.  At the end $\code{acc}\equiv R\beta^{(q-1)/2}=R\chi(\beta)=R\chi(b)$.  The residues $R,0,-R$ are distinct, and \code{acc} lies in $(-q,q)$.
\end{proof}

\subsection{The masked loop and the theorem}

\code{chain\_general} runs all $64$ lanes in one loop.  A lane that has stopped is masked out; when all active lanes have the same degrees the slots are renamed
instead of copied (``lockstep''); otherwise the first $\max\deg+1$ entries are copied under a mask.

\begin{lemma}[The masked loop]\label{lem:lanes}\stproved
Assume the conclusions of Proposition~\ref{prop:ranges} and Lemma~\ref{lem:weights}.  At the start of every iteration of the code's loop, every active lane holds the $U$
and $V$ of the one-lane algorithm on its own polynomial after the same number of iterations (left-aligned, zeros up to index $11$, each with one weight) and its
\code{mm}; every inactive lane has stopped with its final \code{mm} and its stop reason recorded.  Every lane stops within $18$ iterations, and the decision is that of
the one-lane algorithm.
\end{lemma}
\begin{proof}
Induction over the iterations.  At the start every lane holds $f$ and $f'R^{-1}$.  In an iteration, with $\code{imax}$ the largest degree among the active lanes: the
remainder is computed for indices below \code{imax}, and for a lane of smaller degree the extra entries are reductions of $0$, hence $0$, and indices from \code{imax}
on are set to $0$; shifting out leading zeros moves only the lanes whose leading entry is $\equiv0$, and ends with $w=\deg W$ or $W=0$; the zero-remainder,
exponent and constant steps act under the mask of active lanes.  The lockstep branch runs only when every lane is active with the degrees of lane $1$, so renaming
slots performs the same exchange in every lane; in the masked branch all entries beyond \code{imax} are zero, so copying indices up to \code{imax} copies whole
polynomials.  No write reaches an inactive lane's \code{mm}, degrees, flags or $U$, $V$.  By Proposition~\ref{prop:invariant} every lane stops within
$2n-2=18$ iterations, so the code's limit of $40$ is never reached.
\end{proof}

\begin{thmD}[the integer program is exact]\stproved
Assume (S1)--(S3), and let the primes be those \code{pdf\_init} chooses.  For every input buffer, block size and routing state, at every stage and in every lane, the
code's \code{ok} and keep decisions are those of the exact algorithms over $\Fq$ on $f\bmod q$.  So Theorem~C holds for the program as written.  The filter executes no
floating-point operation; its result depends neither on the rounding mode nor on x87 or SSE code nor on \code{-ffast-math}.
\end{thmD}
\begin{proof}
The lanes start with $f\bmod q$ (Lemma~\ref{lem:ldres}); every value is in range (Proposition~\ref{prop:ranges}); the branches and symbols are the exact ones
(Lemma~\ref{lem:weights}); the $-1$ test is exact (Lemma~\ref{lem:eulermont}).  In \code{chain\_block} every statement is an elementwise array operation without
lane-dependent control flow, so each lane runs the algorithm; in \code{chain\_general} Lemma~\ref{lem:lanes} gives the same.  The module has no \code{real} variable.
\end{proof}

\subsection{The size limits, and the guard}\label{sec:limits}

The single-reduction arguments stop at $q\le2^{13}$ (small) and $q\le2^{30}+1$ (large).  Following the actual bounds of the values through the chains gives larger
limits for the whole filter.  Put $\rho(X)=\lfloor(X+qR/2)/R\rfloor$, the bound of Lemma~\ref{lem:montgen} for the result of a reduction of an argument $|x|\le X$.

\begin{lemma}[Bound propagation]\label{lem:propagation}\stproved
Fix a stage with an odd prime $q>5$.  Define: $L=\lfloor(2(2^{21}-1)(q-1)+2^{21}(q-1)+2^{31}q)/2^{32}\rfloor$ for the loading; for \code{chain\_block},
$A_k=L$, $A_{10}=1$, $B_{k-1}=\rho(kA_k)$, and for $m=9,\dots,1$ with $\lambda=B_m$, $\tau=A_{m+1}$: $R_0=\rho(\lambda A_0)$, $R_k=\rho(\lambda A_k)+\rho(\tau B_{k-1})$
(small) or $\rho(\lambda A_k+\tau B_{k-1})$ (large), $P_k=\rho(\lambda R_k)+\rho(R_mB_k)$ (small) or $\rho(\lambda R_k+R_mB_k)$ (large), then $A:=B$, $B:=P$; Euler's
loop with the arguments $B_0^2$, $(q-1)B_0$, $(q-1)^2$; for \code{chain\_general}, $b=\max(L,1)$, $v=\rho(10L)$, then $18$ times $w=2\rho(vb)$ (small) or
$\rho(2bv)$ (large), the argument $(q-1)\max(b,v,w)$, and $b:=v:=\max(b,v,w)$; Euler's loop with $(q-1)^2$.  Call each argument of $\rho$ a \emph{site}, and each
sum of two values of $\rho$ at a small stage a \emph{difference site}.  If every site is below $qR/2$ and every difference site at most $32767$, then every argument
of a reduction is bounded by its site, every result lies in $(-q,q)$, every difference fits $16$ bits, and Lemmas~\ref{lem:weights}, \ref{lem:eulermont} and
Theorem D hold at this stage.
\end{lemma}
\begin{proof}
Induction over the statements, as in Proposition~\ref{prop:ranges}, with the intervals replaced by these numbers: a reduction of an argument with $|x|\le X<qR/2$
returns $|u|\le\rho(X)<q$, and $\rho$ is increasing.  In \code{chain\_general} every slot of every lane is bounded by the current $b$ after each iteration, including
lanes that have stopped; by induction each active lane follows the algorithm and is active for at most $18$ iterations.
\end{proof}

\begin{theorem}[The limits of the whole filter]\label{thm:limits}\stnum
The conditions of Lemma~\ref{lem:propagation} hold for every odd $q$ with $5<q\le8385$ at a small stage and every odd $q$ with $5<q\le1{,}097{,}010{,}265$ at a large
stage, and fail for every odd $q$ in $[8387,32767]$ and in $[1{,}097{,}010{,}267,\ 1{,}100{,}000{,}001]$, first at the $18$th iteration of \code{chain\_general}.  So
Theorem D holds for every choice of small primes $q\equiv1\pmod4$ with $10<q\le8385$ and large primes $q\equiv1\pmod4$ with $10<q\le1{,}097{,}010{,}265$.
\end{theorem}
The computation evaluates the recursion in $128$-bit integers for every odd $q\le32767$ and every odd $q\le1{,}100{,}000{,}001$, and is repeated independently for
every odd $q\le9001$ and around the large limit.

\paragraph{The guard.}  \code{pdf\_init} chooses a usable prime only inside $[11,8385]$ (small) or $[11,1{,}097{,}010{,}265]$ (large), and writes a line on standard error
for any it skips.  The shipped primes are below $2^{13}$ and $2^{25}$, far inside these limits; the guard ensures that no future change of the prime search can
leave the proved region silently.

\paragraph{How sharp the limits are} \stnum.  The individual reductions fail earlier than the whole filter: \code{mont} with operands in $(-2q,2q)$ first returns a
value outside $(-q,q)$ at $q=8241$, and \code{sredc} on a difference $ab-cd$ at $q=2^{30}+23279$ (complete searches); the chains never form those operands.  Above the
filter limits, a directed search on a model of the chains, at every prime $\equiv1\pmod4$ from $8389$ to $19961$ and at $64$ large primes, finds the first results
outside $(-q,q)$ at $15901$ (small) and $1{,}749{,}999{,}961$ (large), and the first wrong verdicts at $18457$ and $2{,}032{,}588{,}133$; those found at primes
\code{pdf\_init} could choose make the real filter fail as well.  So the proved limits are at least $0.53$ and $0.63$ of the true thresholds for such failures, and
probably not sharp; a search cannot establish where the thresholds lie.

\subsection{The start-up check}

\begin{proposition}[Start-up check]\label{prop:kat}\stproved
At the end of \code{pdf\_init} the function \code{arith\_ok} runs the compiled arithmetic on known inputs: \code{mont} and \code{sredc} on $64$ inputs per prime
including the extremes, the relations $qq'\equiv1$, the split of $64$ coefficients including $-2^{63}$ and $2^{63}-1$, and the loads of the filter themselves on a block of
$64$ polynomials with these coefficients.  The expected values are computed with $+$, $-$, $\times$, \code{modulo} and exact division only, which do not depend on (S2)
or (S3).  If any comparison fails, \code{pdf\_init} writes one line on standard error and chooses no prime; then every flag is $1$ and the exact test alone decides.
\end{proposition}
\begin{proof}
With no prime chosen the stage loop does not run, and every flag is set to $1$.  The rest is what the comparisons test.
\end{proof}

\paragraph{What the check is for.}  It does not guard the run-time environment, as a check of a floating-point filter would have to: integer instructions have no
rounding mode, so a given binary computes the same values on every processor of its architecture.  What (S2) and (S3) depend on is the compiler that built the
binary, and that is fixed when the program is built.  The release builds use several: the MinGW-w64 cross compiler for Windows, gfortran for x86-64 and aarch64
Linux, Homebrew's gfortran on macOS (arm64), the Android cross compiler (aarch64) and FreeBSD's gfortran, while the tests of Section~\ref{sec:evidence} ran on x86-64
with GCC~15.2 and 16.2 only.  For the other builds the start-up check is the one test of their compiled integer code, and for a cross-compiled build the target
machine is the first place where that code can run at all.  The check costs no measurable time and fails safe: the filter switches itself off and the exact test
decides every polynomial.  It is not a proof.  It runs the compiled loads exactly as the filter does, so a fault in the split (such as the top piece shifted
logically) fails it, but not the copies of the reductions that the compiler inlines into the chains; (S2) and (S3) therefore remain named assumptions
(Section~\ref{sec:ledger}).

\paragraph{Instruction sets.}  The application contains an AVX2 build and a baseline (SSE2) build of the filter and uses the AVX2 build when the CPU reports AVX2 and
FMA.  The AVX2 build is compiled with \code{-march=x86-64 -mavx2 -mno-popcnt}, so it can contain nothing the run-time test does not check: its $114$ distinct
instructions are base x86-64 instructions and VEX-encoded AVX and AVX2 instructions; the baseline object uses SSE and SSE2 only \stnum.

%======================================================================
\section{The GPU filter}\label{sec:gpu}

The GPU kernel runs the enumeration's innermost loops and, for every polynomial it counts, the one-lane \code{chain\_general} at the primes of its table.  A
polynomial it rejects is never shown to the CPU, so the kernel's filter must be sound by itself; its survivors go through the CPU filter and the exact test.
It uses unsigned $32$-bit Montgomery arithmetic with $R=2^{16}$ and residues in $[0,2q)$ (zero modulo $q$ exactly when $0$ or $q$), which needs no
multiply-high instruction and no comparison per reduction.  The table row of $q$ holds $q$, the loading weights $2^{16j}\bmod q$, $(q-1)/2$, $q^-=-q^{-1}\bmod2^{32}$
and $r^-=(-R)\bmod q$.

\begin{lemma}[\code{redc}]\label{lem:redc}\stproved
Let $q<8192$ be odd, $0\le t<8q^2$ and $m=tq^-\bmod2^{16}$.  Then $t+mq<2^{30}$ is divisible by $2^{16}$, and $r=(t+mq)/2^{16}\equiv t\,2^{-16}\pmod q$ with
$0\le r<2q$.
\end{lemma}
\begin{proof}
$mq\equiv-t\pmod{2^{16}}$; $t+mq<8q^2+2^{16}q<2^{30}$; $r<t/2^{16}+q<q(8q/2^{16}+1)<2q$.  Unsigned arithmetic in OpenCL C and CUDA wraps modulo $2^{32}$ by
definition, and none of these values wraps.
\end{proof}

\begin{thmE}[the GPU filter is exact]\stproved
For a prime $q\equiv1\pmod4$ with $11\le q<8192$, its table row, and any $64$-bit coefficients $c_1,\dots,c_{10}$, the kernel rejects $f=x^{10}+c_1x^9+\dots+c_{10}$
at $q$ exactly when $\leg{\disc f}{q}=-1$.  The table's primes are usable for $S$, so a polynomial the kernel rejects does not pass.
\end{thmE}
\begin{proof}
\emph{Loading.}  With the $16$-bit limbs $u_j$ of $x\bmod2^{64}$, the integer sum $\sum_ju_j(2^{16j}\bmod q)$ is at most $(2^{16}-1)(1+3(q-1))<2^{31}$; subtracting
$2^{64}\bmod q$ for $x<0$ and reducing gives $x\bmod q$.  \emph{Ranges.}  Every residue is in $[0,2q)$, and every argument of \code{redc} is below $8q^2$:
$(10-i)a_i<10q$, $l_vU(i)+(2q-l_u)V(i)<(2q-1)^2+2q(2q-1)<8q^2$, products of two residues below $4q^2$.  As $2q-l_u\equiv-l_u$, each new coefficient is the exact
$W(i)$ times $R^{-1}$.  \emph{The chain} is the one-lane algorithm of Section~\ref{sec:chains} without masks; the zero test \code{w[0]==0 || w[0]==q} is exact on
$[0,2q)$, entries beyond a degree are $0$, and it stops within $18<40$ iterations.  As in Lemma~\ref{lem:weights}, every polynomial is its exact counterpart times a
power of $R^{-1}$, so the branches are exact and $\chi(\code{mm})$ is the exact symbol, which is $\leg{\disc f}{q}$ by Theorem B.  \emph{Euler's criterion} runs as in
Lemma~\ref{lem:eulermont} with $\code{acc}\in[0,2q)$, so $\chi=-1$ exactly when $\code{acc}\in\{r^-,r^-+q\}$.
\end{proof}

\paragraph{The table.}  The table takes the smallest usable primes in $(256,8192)$, at most $24$.  Every set $S$ of one to three primes below $60$ (the application
allows $|S|\le3$) has at least $24$ there \stnum, while a range below $1024$ would leave fewer than the $8$ that GPU mode requires for $592$ of these $833$ sets.  The
smallest primes come first because Euler's criterion takes one step per bit of $(q-1)/2$.

\paragraph{Host and device code must match.}  Theorem E needs the kernel to read the table the program builds.  For OpenCL this holds by construction: the kernel
source is part of the program.  The CUDA kernel is a fat binary embedded at build time and the MUSA kernel a library loaded at run time, so either could be left over
from an older kernel that reads the table differently and rejects passing polynomials.  The table layout therefore carries a number (format~$2$), the kernel's entry
point carries it in its name (\code{nf\_enum\_t2}), and the MUSA library reports it and checks that every table row has the integer layout.  A CUDA program with
an old fat binary stops at start-up; a MUSA program refuses an old library and runs on the CPU; an old program with a new library stops \stnum.  After any change of
the kernel the CUDA fat binary and the MUSA library must be rebuilt.

\paragraph{Run-time checks on the device.}  The kernel was tested on one card, an NVIDIA GTX~1650 (Section~\ref{sec:evidence}).  Under OpenCL its source is
compiled at start-up by the graphics driver of each machine, and the CUDA and MUSA kernels run on cards the tests did not use, so the code a device runs has in
general not been tested.  Unlike the CPU's start-up check (Proposition~\ref{prop:kat}), which stands in for tests of the release builds, these checks remain
necessary however many builds are tested.  At start-up every back end runs a self-test kernel of the same source once: \code{redc}, the loading, the verdict at
every table prime for $240$ polynomials and the full filter, $7{,}920$ values for $S=\{2,5\}$, which the host computes without the kernel's code (exact arithmetic
modulo $q$, the symbol from the resultant).  On a difference, or with device code that has no self-test, the task runs on the CPU.  In the window at the start of a
task, where CPU loops run next to the GPU, the program compares polynomial counts; with \code{NF\_GPUCHECK=1} the CPU also applies the kernel's filter, compiled for
the CPU, to the whole task, and the GPU's survivors must be the same, in the same order, or the task stops (about a microsecond per polynomial).

%======================================================================
\section{What the exact test of GetDecics computes}\label{sec:poldisc}

In the application the filter's rejections take the place of verdicts of GetDecics' own test, \code{polDiscTest\_cpu}, so it matters what that test computes.
It follows Cohen's sub-resultant algorithm \cite[Algorithm~3.3.7]{cohen} on $A=f$ and $B=f'$: with $g=h=1$, it repeats $\delta=\deg A-\deg B$,
$R=\mathrm{prem}(A,B)$ (the pseudo-remainder $\lc(B)^{\delta+1}A\bmod B$), $A\leftarrow B$, $B\leftarrow R/(gh^\delta)$, $g\leftarrow\lc A$,
$h\leftarrow h^{1-\delta}g^\delta$, until $\deg B\le0$.  The numbers grow to about a thousand bits, so the test uses GMP; to save time it does the first two
iterations in $64$-bit integers (\code{int64}) and only then switches to GMP.  At the end it removes every factor of the primes of $S$ from its result \code{polDisc}
and asks whether what remains is a perfect square.

\begin{lemma}[Sub-resultants]\stquoted
In this algorithm every division $R/(gh^\delta)$ and $g^\delta/h^{\delta-1}$ is exact in $\Z[x]$ resp.\ $\Z$, every $g,h$ is nonzero, and the final
$s\,t\,h^{1-\deg A}\lc(B)^{\deg A}$ is $\Res(A_0,B_0)$, where $t$ is the product of the contents' powers and $s$ a sign \cite[Section~3.3]{cohen}.
\end{lemma}

\begin{proposition}[What the exact test computes]\label{prop:poldisc}\stproved
Let $f$ be monic of degree $10$ with \code{int64} coefficients, and suppose (H): no \code{int64} value computed in the first two iterations of
\code{polDiscTest\_cpu} leaves the \code{int64} range.  Then $|\disc f|=r^2|\code{polDisc}|$ for a nonzero rational $r$, and \code{polDiscTest\_cpu} sets the flag
to $1$ exactly when $f$ passes.
\end{proposition}
\begin{proof}
By (H) every \code{int64} expression has its mathematical value and the conversions to GMP are exact, so the computation is one over $\Z$.
\emph{Initialisation:} $\lc(f')=10$, so $\cont(f')$ divides $10$; the coefficients of the odd powers of $f'$ are even, so $2\mid\cont(f')$ exactly when the even-power
coefficients are even, which the code tests, and it tests $5$ likewise; it divides $B$ by $\cont(f')$, and $\cont(f)=1$.
\emph{Pseudo-remainders:} each elimination step multiplies by $\lc B$ and subtracts a multiple of $B$ that removes the leading term, $\delta+1$ times in all
(counting steps whose leading term is already zero as multiplications by $\lc B$); the result has degree below $\deg B$ and differs from $\lc(B)^{\delta+1}A$ by a
multiple of $B$, so it is $\mathrm{prem}(A,B)$ by uniqueness of division in $\Q[x]$.
\emph{Reduction:} the code's scale factors are $gh^\delta$ and its update of $h$ is Cohen's; by the lemma the divisions are exact, so C's \code{/} and
\code{mpz\_divexact} return the exact quotients.
\emph{The end:} if the final constant $b_0$ is $0$ the resultant is $0$ and the flag is $0$.  Otherwise, with $d=\deg A$, $\Res(f,f')=s\,t\,h^{1-d}b_0^d$ with
$t=\cont(f')^{10}$, and the code takes $\code{polDisc}=h$ for even $d$ and $b_0$ for odd $d$; the quotient is $(b_0/h)^d$ resp.\ $(b_0/h)^{d-1}$, an even power.  With
Proposition~\ref{prop:F1}, $|\disc f|=\cont(f')^{10}|h^{1-d}b_0^d|=r^2|\code{polDisc}|$.
\emph{The square test:} the code removes every factor of each $p\in S$ from $|\code{polDisc}|$ and tests the rest with \code{mpz\_perfect\_square\_p}; that holds
exactly when $v_p(\code{polDisc})$ is even for all $p\notin S$, which the factor $r^2$ does not change.  This is Definition~\ref{def:pass}.
\end{proof}

\paragraph{Hypothesis (H).}  The two \code{int64} iterations multiply coefficients of $f$ by coefficients of $f'$, and then by coefficients of the first remainder;
the code does not check that these products fit into $63$ bits and relies on the polynomials it is given.  The polynomials of the enumeration are bounded by its
search region, and the author of GetDecics states that polynomials generated for the fields under test do not overflow these steps \streport.  The source shows the
same design: it sends one branch of the second iteration, where ``the result would exceed a long'', to GMP.  For a polynomial outside the region the overflow is
undefined behaviour in C as GetDecics is built; Passiflora compiles the test with \code{-fwrapv}, so that the value wraps around modulo $2^{64}$ by definition (as it
does in practice anyway).  The GMP steps that follow then work on a wrong number, and the verdict can be wrong in either direction.

\paragraph{Measurement.}  The program \code{test\_poldisc.cpp} links the compiled \code{polDiscTest\_cpu} of GetDecics, compares its flag with the definition computed
independently by PARI (the discriminant, the primes of $S$ removed, the square test), and replays the \code{int64} steps in $128$-bit integers to detect an overflow.
It runs the real candidates of the application, the results of the test work units, and families built to test the proposition (random coefficients,
polynomials built to pass, square factors), which the enumeration does not produce (Table~\ref{tab:poldisc}).

\begin{table}[ht]
\centering\footnotesize\setlength{\tabcolsep}{4pt}
\caption{The compiled \code{polDiscTest\_cpu} of GetDecics against the definition (PARI), and the Passiflora filter, on $10{,}735{,}681$ polynomials ($S=\{2,5\}$).
``ovf'': an \code{int64} step of GetDecics overflows; ``differ'': the flag of GetDecics is not the definition's.}\label{tab:poldisc}
\begin{tabular}{@{}lrrrrrrrr@{}}
\toprule
& & \multicolumn{2}{c}{pass} & & \multicolumn{2}{c}{differ} & \multicolumn{2}{c}{filter} \\ \cmidrule(lr){3-4} \cmidrule(lr){6-7} \cmidrule(l){8-9}
family & polynomials & GetDecics & definition & ovf & no ovf & ovf & rejects & rejects passing \\ \midrule
random, $|a_i|\le3$ & 3,000,000 & 458 & 458 & 0 & 0 & 0 & 2,930,945 & 0 \\
random, $|a_i|\le100$ & 2,000,000 & 0 & 0 & 0 & 0 & 0 & 1,999,948 & 0 \\
random, $|a_i|\le10^6$ & 1,000,000 & 0 & 0 & 999,961 & 0 & 0 & 1,000,000 & 0 \\
built to pass, small & 2,000,000 & 1,999,939 & 1,999,939 & 0 & 0 & 0 & 0 & 0 \\
built to pass, large & 1,000,000 & 155,735 & 1,000,000 & 844,277 & 0 & 844,265 & 0 & 0 \\
square factor & 1,000,000 & 0 & 0 & 0 & 0 & 0 & 0 & 0 \\
real candidates & 733,855 & 0 & 0 & 0 & 0 & 0 & 733,855 & 0 \\
the $83$ results, transformed & 1,826 & 1,826 & 1,826 & 0 & 0 & 0 & 0 & 0 \\ \midrule
all & 10,735,681 & 2,157,958 & 3,002,223 & 1,844,238 & 0 & 844,265 & 6,664,748 & 0 \\
\bottomrule
\end{tabular}
\end{table}

\paragraph{What it shows} \stnum.  Wherever (H) holds, GetDecics agrees with the definition: on all $8{,}891{,}443$ polynomials without an overflow, as the
proposition says.  Outside (H), with coefficients far outside the search region, the \code{int64} steps overflow and the answer is no longer the definition's.
Coefficient size alone is not the measure: the real candidates have coefficients up to $2.1\cdot10^9$ and never overflow, because their coefficients grow with the
degree as those of polynomials with roots of moderate size do, which is what the search region produces.  What matters is the largest value the \code{int64} steps
form.

\paragraph{The evidence on real data} \stnum.  On the $146{,}771$ candidates captured from seven real work units the largest \code{int64} value is $2^{59.45}$, a factor
of about $12$ below the limit $2^{63}$; on the $83$ results of the test work units it is $2^{48.5}$; on all $733{,}855$ real candidates and shifts of
Table~\ref{tab:poldisc} no step overflows.  A copy of the application that replays the \code{int64} steps of every polynomial before the filter (the $128$-bit replay of
\code{test\_poldisc.cpp}) found no overflow among the $14{,}846{,}153{,}955$ polynomials of the $31$ test work units, all of which gave their truth files.  This
supports the author's statement, but it is a measurement, not a proof from the search bounds: (H) remains a named assumption, with three and a half bits of measured
headroom.

\paragraph{The mechanism, illustrated.}  The following polynomial is constructed; the enumeration does not produce it.  One of the $83$ results in the truth files is
$f=x^{10}+600x^7+1350x^6-800x^5+115000x^4+405000x^3+265625x^2-540000x+185000$.  Multiplying its roots by $3$ gives
\begin{multline*}
g=3^{10}f(x/3)=x^{10}+16200x^7+109350x^6-194400x^5+83835000x^4+885735000x^3\\
+1742765625x^2-10628820000x+10924065000 ,
\end{multline*}
which defines the same field, and $\disc g=3^{90}\disc f$, so $g$ passes as $f$ does.  Its largest \code{int64} value would be about $2^{67}$, and GetDecics
rejects it; with the roots multiplied by $2$ the largest value is about $2^{56.5}$ and GetDecics accepts it.

\paragraph{Passiflora and (H).}  Passiflora keeps the exact test as GetDecics has it (its patch changes only the test's conversions for a $32$-bit \code{long} and adds
comments, and \code{-fwrapv} only defines what an overflow does), so that its output is the stock program's.  The filter rejects only polynomials that provably do not pass, so it never needs (H) for its own soundness;
(H) enters only where the filter's verdicts are compared with those of the exact test (contract (C1)).  Outside (H) the filter can differ from the stock program only
by rejecting a polynomial that the exact test wrongly accepts.  That did not happen in the tests (last column of Table~\ref{tab:poldisc}), but it can be constructed
\stnum.  When $B_9=10/\cont(f')$ is even, the \code{int64} steps see the constant term $a_0$ only through $B_9^2a_0$, and $4\mid B_9^2$, so adding $2^{62}$ to $a_0$
leaves every wrapped value unchanged: $h=x^{10}+x^9-3x^8-2x^7+x^6-x^5-x^3+3x^2+2x-2$ passes, and GetDecics also passes $h+2^{62}$, which the filter rejects.  The
polynomials written out would be the same (PARI's field-discriminant test rejects $h+2^{62}$), but the counters would differ.  Such constant terms are far outside any
search region.  A check of (H) on every polynomial would cost about half of the filter's time; a checked version of the two iterations with a fallback to GMP would
cost the stock program about $1\%$ per exact test and is the change to propose to GetDecics, so that both programs stay identical.

%======================================================================
\section{The contract with the application}\label{sec:contract}

\begin{description}
\item[(C1) Soundness.] A flag is $0$ only if the \code{polDiscTest\_cpu} of GetDecics would reject the polynomial.
\item[(C2) Safety.] No integer overflow, out-of-bounds access, read of an undefined value, other undefined behaviour, or illegal instruction.
\item[(C3) Integration.] The flags reach the exact test as its \code{polyMask} argument expects, and the $31$ test work units give their \code{.truth} files.
\end{description}

\begin{proposition}[The interface]\label{prop:contract}\stproved
After \code{pdf\_init(numP, pSet)}, \code{pdf\_filter(pol, n, flag)} with $n\le0$ reads and writes nothing; with $n\ge1$ it reads only the $n$ rows of $11$ coefficients,
writes every flag with $0$ or $1$, and writes $0$ only for a polynomial whose leading coefficient is $1$ and which does not pass for $S=\{\code{pSet}\}$.  Before the
first \code{pdf\_init}, and after a failed start-up check, every flag is $1$.  The module state left by earlier calls changes the routing, never a flag.
\end{proposition}
\begin{proof}
The loads read only rows whose indices are in the work lists, which hold $1,\dots,n$; every flag is written.  A flag $0$ means a rejection at a chosen prime, which by
Theorem D and Proposition~\ref{prop:rule} means that $f$ does not pass.  With no prime chosen every flag is $1$ (Proposition~\ref{prop:kat}).  Theorem D holds for
every routing state, and the work arrays are reallocated when they are too small.
\end{proof}

\begin{proposition}[The glue]\label{prop:glue}\stproved
The application calls the filter into its flag array and then \code{polDiscTest\_cpu} with the same array as \code{polyMask}.  Afterwards the flag is $0$ where the
filter wrote $0$ and the exact test's own flag elsewhere; under (C1) that is the exact test's flag everywhere, the output of the stock program.
\end{proposition}
\begin{proof}
\code{polDiscTest\_cpu} reads the mask at index $k$ before it writes the flag at index $k$, touches no other index in that iteration, and writes a flag on every path.
Every polynomial of the enumeration is monic (the enumeration sets the leading coefficient to $1$ once), and the filter and the exact test receive the same $S$.
\end{proof}

With Proposition~\ref{prop:poldisc}, these give (C1) under (H).

\paragraph{Safety} \stproved.  Overflow: the range proofs; the call counter is kept modulo $64$.  Bounds: every index is at most $n$ or the lane count, and the chains'
indices stay in $0..11$ (Lemma~\ref{lem:lanes}).  Undefined values: padding lanes hold $x^{10}$; a flag that \code{chain\_general} does not set is set before it is read.
Illegal instructions: the AVX2 object runs only when the CPU reports AVX2 and FMA, and contains nothing beyond AVX2.  Language-level undefined behaviour: none apart from
the conversions and shifts of (S2), (S3), which the start-up check tests.

\paragraph{Further conditions.}  (A1) At most three primes in $S$ \stproved: GetDecics keeps $S$ in \code{int pSet[3]}, and so do the patch's copies; the program refuses a
work unit with more primes right after it reads $S$, before any array is filled (tested with four and five primes; all $31$ work units have $|S|=2$).  Named
assumptions: (A2) one thread: the filter keeps module state, and the application calls it only from the thread that runs the enumeration; (A3) hypothesis (H) of
Proposition~\ref{prop:poldisc}.

%======================================================================
\section{Evidence}\label{sec:evidence}

The theorems are about the program text and the mathematics; the compiler, the device and the code of GetDecics are trusted only as far as the following checks
show.  CPU runs are on Linux (WSL2, GCC and gfortran 15.2) unless stated.

\paragraph{The compiled filter on complete sets at tiny primes} \stnum.  The filter was built from the shipped source with its prime search moved to one tiny prime
$q\in\{13,17,29,37,41\}$ (four constants changed), once in the $16$-bit and once in the $32$-bit stage, and every polynomial of a set was sent through the AVX2 and the
baseline build, each with the normal routing and with every lane forced through \code{chain\_general}: eight runs per polynomial.  Every flag and every rejection stage was
compared with an exact oracle (Euclid with inverses).  The sets: all $13^9$ depressed monic polynomials modulo $13$ (every monic polynomial modulo $13$ is a shift of one
of them); at $17$ all with $a_9=a_8=0$, at $29$, $37$, $41$ all with $a_9=\dots=a_6=0$ (which send every chain into \code{chain\_general} and its lockstep branch); and at
each prime $200{,}000$ random polynomials with all their shifts and coefficients lifted by large multiples of $q$.  Altogether $25{,}518{,}310{,}785$ polynomials and
$204{,}146{,}486{,}280$ verdicts: none wrong.

\paragraph{The arithmetic alone} \stnum.  The reductions, cut out of the source (not retyped) and compiled with the filter's flags and at \code{-O0}, against exact
models: \code{mont} on every pair of operands for every odd prime below $2^{13}$ ($333{,}937{,}622{,}331$ pairs) and every odd $q$ in $[8193,8303]$; \code{sredc} on
windows around its limits and on random inputs; the loading on $101{,}072{,}896$ coefficients for $6{,}169$ moduli; the bound propagation of
Theorem~\ref{thm:limits}.  With GCC~15.2 (Linux) and 16.2 (MSYS2): no error.

\paragraph{An exact oracle, and deliberate faults} \stnum.  A test that only feeds polynomials that must pass cannot see a filter that rejects too little.  The
oracle test computes the symbol at every filter prime independently, builds its inputs to reach every stage and path of the filter (controlling the symbols by the
Chinese remainder theorem), adds metamorphic checks, and reads a test hook that reports the stage of every rejection.  Seven deliberate faults of the integer
arithmetic, each one changed line, against $1{,}000{,}000$ polynomials (total failures, AVX2 build):

\begin{center}\footnotesize
\begin{tabular}{@{}lrr@{}}
\toprule
fault & oracle test & ``must pass'' test \\ \midrule
none & 0 & 0 \\
zero test $\pm q$ in \code{chain\_block\_s} & 360 & 0 \\
zero test $\pm q$ in \code{chain\_general\_s} & 23,283 & 349 \\
Euler: $-R$ recognised only as one representative & 57,624 & 0 \\
top piece of a coefficient shifted logically (start-up check: filter off) & 643,008 missed & 0 \\
wrong loading constants (start-up check: filter off) & 643,008 missed & 0 \\
parity of \code{chain\_general} & 236,849 & 161,538 \\
last stage skipped & 7,280 & 0 \\
\bottomrule
\end{tabular}
\end{center}

The oracle test sees all seven.  The floating-point faults (x87, rounding up, down, toward zero) give $0$ on the integer filter, as Theorem D says, while on the
earlier floating-point filter the oracle test saw $66$ false and $296$ missed rejections under x87 and $37{,}329$ false under rounding up.  Primes moved across the size
limits behave as the proofs predict: no wrong verdict inside them; beyond them wrong verdicts appear, at some primes only for directed inputs.  With $24$ small primes
the generated polynomials reach the large stages only through constructions that control every symbol; the large-stage code is instruction for instruction
(\code{objdump}) that of a build with six small primes below $2^{10}$, which the oracle test ran with $788{,}246$ of $1{,}000{,}000$ polynomials reaching the large
stages and no failure.  The $146{,}771$ candidates captured from seven real work units pass through the filter with the oracle's verdicts.

\paragraph{The exact test} \stnum.  Table~\ref{tab:poldisc} and Section~\ref{sec:poldisc}.

\paragraph{The GPU filter} \stnum.  The kernel's filter, compiled for the CPU from the kernel source itself, against the exact oracle prime by prime:
$6{,}000{,}000$ polynomials at the $24$ table primes and $3{,}000{,}000$ at tiny primes and just below the limit $8192$, $228{,}000{,}000$ checks, no false and no
missed rejection.  On an NVIDIA GTX~1650 (OpenCL, and CUDA with the kernel compiled by NVRTC), the application gave the truth files of $8$ work units through OpenCL
and $4$ through CUDA, with the start-up self-test passing in each; a full comparison of one unit ($3{,}759$ survivors) found no difference, and injected faults and a
kernel without the self-test were caught.  The MUSA path was run with a stand-in runtime.

\paragraph{The work units} \stnum.  The application, built on Linux from GetDecics at \code{638ca4e} with the patch and the repository's compile lines (BOINC replaced
by stand-ins), gave the \code{.truth} files of all $31$ test work units with the AVX2 build and with the baseline build.  With \code{NF\_VERIFY=1} (the exact test also
on every polynomial, stopping at the first difference), five units in both builds ran $20{,}008{,}960$ polynomials each without a difference.

\paragraph{Other sets $S$} \stnum.  The proofs hold for every $S$.  The tests were also run for $\{2,11\}$, $\{5,11\}$, $\{2,5,11\}$, $\{2,3,5\}$ and $\{3,7\}$ without a
failure, on fewer polynomials than for $\{2,5\}$.

%======================================================================
\section{Performance}\label{sec:speed}

\paragraph{CPU} \stnum.  Table~\ref{tab:speed} gives the application's own timing split (\code{NF\_PROFILE}) over the first $60{,}006{,}400$ polynomials of three
work units, seven runs per filter, the filters alternating, one run at a time on an otherwise idle machine.  The programs differ only in the filter: the earlier
floating-point filter; an integer filter with Barrett reduction ($32$-bit lanes for the small primes, $64$-bit products for the large ones); and the Montgomery filter
in three steps, each with one change more than the step before.  On \code{sf5} and \code{sf3} almost every polynomial takes \code{chain\_block}; on \code{sf4} every
polynomial takes \code{chain\_general} at every prime.

\begin{table}[ht]
\centering\footnotesize\setlength{\tabcolsep}{4pt}
\caption{CPU filter plus exact test (the application's polDisc time) in seconds for the first $60{,}006{,}400$ polynomials, median of seven runs; in brackets the
change against the floating-point filter.  Montgomery: $6+24$ primes (small $+$ large) with scalar loads; the same primes with vectorised loads; then $24+6$ primes
with these loads, the filter described here.  Work units \code{sf5\_DS-14x11\_Grp33043of705883}, \code{sf3\_DS-16x270\_Grp74of3932160},
\code{sf4\_DS-13x271-25\_Grp98of200000}.}\label{tab:speed}
\begin{tabular}{@{}llrrrrr@{}}
\toprule
& & & & \multicolumn{3}{c}{Montgomery} \\ \cmidrule(l){5-7}
& & & & scalar loads & \multicolumn{2}{c}{vectorised loads} \\ \cmidrule(lr){5-5}\cmidrule(l){6-7}
work unit & build & floating point & Barrett & $6+24$ primes & $6+24$ primes & $24+6$ primes \\ \midrule
\code{sf5} & AVX2     & $5.84$  & $12.03$ ($+106\%$) & $6.65$ ($+14\%$)  & $5.49$ ($-6\%$)  & $5.03$ ($-14\%$) \\
           & baseline & $8.86$  & $24.47$ ($+176\%$) & $8.80$ ($-1\%$)   & $7.96$ ($-10\%$) & $7.21$ ($-19\%$) \\
\code{sf3} & AVX2     & $5.44$  & $11.87$ ($+118\%$) & $6.52$ ($+20\%$)  & $5.65$ ($+4\%$)  & $5.23$ ($-4\%$) \\
           & baseline & $8.20$  & $22.52$ ($+175\%$) & $7.89$ ($-4\%$)   & $7.33$ ($-11\%$) & $6.59$ ($-20\%$) \\
\code{sf4} & AVX2     & $13.37$ & $16.09$ ($+20\%$)  & $8.55$ ($-36\%$)  & $7.64$ ($-43\%$) & $6.99$ ($-48\%$) \\
           & baseline & $21.88$ & $32.81$ ($+50\%$)  & $11.25$ ($-49\%$) & $10.86$ ($-50\%$) & $9.28$ ($-58\%$) \\ \bottomrule
\end{tabular}
\end{table}

\emph{Loading.}  With scalar loads the small-prime stages gathered three arrays of $21$-bit pieces through the block's indices, a loop the compiler does not
vectorise; cycle counters in the application showed these loads ($114$ cycles per polynomial on \code{sf5}, against $61$ in the floating-point filter) to account for
the whole gap with AVX2.  The vectorised loads copy each polynomial's ten coefficients into a block buffer and split and reduce them on contiguous lanes ($44$ cycles).
The arithmetic and the compiled chains are the same in both.  \emph{The prime split.}  With $24$ small primes almost no polynomial reaches the large stages, which cost
about twice as much as in the floating-point filter ($0.4$ instead of $30.7$ cycles per polynomial there, $6.2$ more in the small stages, the filter $8.6\%$ faster).
The resulting filter is faster than the floating-point one in every row: by $14\%$ and $4\%$ with AVX2 on \code{sf5} and \code{sf3} (the $4\%$ is less than the
spread of the runs, $5$--$31\%$ between the fastest and the slowest run of one filter), by $19$--$20\%$ with the baseline build, and by $48$--$58\%$ on \code{sf4}.
Barrett reduction is the slowest throughout: AVX2 has no $64$-bit vector multiplication, and a $32$-bit lane holds half as many residues per register as the $16$-bit
lanes of the Montgomery stage.  Timings of different series differ by up to about $10\%$, so only filters of one series are compared.

\emph{Why $24+6$} \stnum.  A cost model, measured stage by stage with cycle counters in the application (five work units, ten splits, both builds, five runs each),
predicts the filter time of every split to within $1\%$.  Each stage keeps half of the polynomials whatever the size of its prime ($0.5001$ at the small primes,
$0.4993$--$0.5007$ at the large ones), a call stops as soon as it has no polynomial left, a large stage costs $3.1$--$4.2$ times a small one per block, and the exact
test costs $6.7$--$8.0\,\mu$s per kept polynomial.  An exchange argument shows that small primes belong first \stproved; of the splits of $30$ stages, $30+0$ is the
cheapest and $24+6$ is within $0.14\%$ of it, and $30$ stages are within $0.10\%$ of the best number.  All three prime sets of the project's searches ($\{2,5\}$,
$\{2,11\}$, $\{5,11\}$) have more than $100$ usable primes below $2^{13}$, and the first stage alone takes almost half of the filter's time.

\paragraph{GPU} \stnum.  Table~\ref{tab:gpuspeed} gives one series of the kernel in four versions: the floating-point kernel, and the integer kernel in three
steps, each with one change more than the step before (the integer arithmetic with $14$ primes below $2^{10}$; the table of $24$ primes below $2^{13}$; the start-up
self-test).  The versions alternate, and every run gave its truth file.  The integer kernel takes $9$--$12\%$ less device time than the floating-point one; the wider
table adds $0.3$--$0.9\%$; the self-test, which runs once at start-up, leaves the device time unchanged.  A whole run takes $3$--$10\%$ less on the three longer
units; on \code{sf5}, where start-up and the window check take most of the run, the differences are inside the spread of the runs.  An integer kernel with
$R=2^{32}$, canonical residues and \code{mul\_hi} was also exact but $47\%$ slower than the floating-point kernel.

\emph{Measurement conditions.}  The test machine is a laptop (Intel i5-12500H, GTX~1650) on mains power.  CPU load lengthens the device time on this machine:
with all $16$ hardware threads busy the device time on \code{sf4} rose from $2651$ to $3123$\,ms ($+18\%$) and returned to $2666$\,ms afterwards, on \code{sf7}
from $6975$ to $7015$\,ms.  The GPU's core clock stayed at its maximum, $1785$\,MHz, with no throttle reason; the GPU was idle for longer stretches between
launches and its memory clock dropped to its idle value at times, a likely cause that was not established.  The GPU times of Table~\ref{tab:gpuspeed} were
therefore taken with the CPU otherwise idle.

\begin{table}[ht]
\centering\footnotesize\setlength{\tabcolsep}{4pt}
\caption{GPU filter on a GTX~1650 (CPU otherwise idle, GPU core clock at its maximum $1785$\,MHz whenever busy), one series, medians; in brackets the change against
the floating-point kernel.  ``device'': the device time of the filter kernel's launches (OpenCL, a build instrumented for this measurement, fixed launch size,
window check off), $7$ runs, spread $0.5$--$2.5\%$.  ``OpenCL'', ``CUDA'': the whole run with the program's defaults, $5$ runs, spread up to $13\%$ on \code{sf5}, up to
$8\%$ on the others.  Integer kernel: for $S=\{2,5\}$ the $14$ usable primes in $(256,1024)$; the $24$ smallest usable primes below $8192$; then with the start-up
self-test, the program described here.  Work units as in Table~\ref{tab:speed} and \code{sf7\_DS-12x12\_Grp881552of1000000}.}\label{tab:gpuspeed}
\begin{tabular}{@{}llrrrr@{}}
\toprule
& & & \multicolumn{3}{c}{integer, Montgomery $R=2^{16}$} \\ \cmidrule(l){4-6}
& & & $14$ primes below $2^{10}$ & \multicolumn{2}{c}{$24$ primes below $2^{13}$} \\ \cmidrule(lr){4-4}\cmidrule(l){5-6}
work unit & time & floating point & no self-test & no self-test & self-test \\ \midrule
\code{sf5} & device (ms) & $1191$  & $1044$ ($-12.3\%$)  & $1054$ ($-11.5\%$)  & $1054$ ($-11.5\%$) \\
           & OpenCL (s)  & $2.68$  & $2.70$ ($+1\%$)     & $2.67$ ($0\%$)      & $2.62$ ($-2\%$) \\
           & CUDA (s)    & $2.51$  & $2.49$ ($-1\%$)     & $2.49$ ($-1\%$)     & $2.46$ ($-2\%$) \\
\code{sf3} & device (ms) & $15029$ & $13592$ ($-9.6\%$)  & $13627$ ($-9.3\%$)  & $13629$ ($-9.3\%$) \\
           & OpenCL (s)  & $16.76$ & $15.41$ ($-8\%$)    & $15.37$ ($-8\%$)    & $15.32$ ($-9\%$) \\
           & CUDA (s)    & $18.04$ & $16.69$ ($-7\%$)    & $16.74$ ($-7\%$)    & $16.77$ ($-7\%$) \\
\code{sf4} & device (ms) & $3027$  & $2662$ ($-12.1\%$)  & $2674$ ($-11.7\%$)  & $2692$ ($-11.1\%$) \\
           & OpenCL (s)  & $7.97$  & $7.55$ ($-5\%$)     & $7.56$ ($-5\%$)     & $7.61$ ($-5\%$) \\
           & CUDA (s)    & $7.86$  & $7.58$ ($-4\%$)     & $7.45$ ($-5\%$)     & $7.64$ ($-3\%$) \\
\code{sf7} & device (ms) & $7902$  & $6989$ ($-11.6\%$)  & $7013$ ($-11.2\%$)  & $7011$ ($-11.3\%$) \\
           & OpenCL (s)  & $9.59$  & $8.63$ ($-10\%$)    & $8.78$ ($-8\%$)     & $8.80$ ($-8\%$) \\
           & CUDA (s)    & $9.52$  & $8.94$ ($-6\%$)     & $8.89$ ($-7\%$)     & $8.99$ ($-6\%$) \\ \bottomrule
\end{tabular}
\end{table}

%======================================================================
\section{What it all rests on}\label{sec:ledger}

\begin{center}\footnotesize
\begin{tabularx}{\textwidth}{@{}>{\raggedright\arraybackslash}p{0.2\textwidth}>{\raggedright\arraybackslash}p{0.12\textwidth}X@{}}
\toprule
claim & status & rests on \\ \midrule
(C1) soundness & \stproved & Theorems A--D, Proposition~\ref{prop:poldisc}, the interface; assumes (H), compiler correctness, (S2), (S3) \\
(C2) safety & \stproved & range proofs, size limits, the masked-loop lemma; assumes (S2), (S3) (start-up check), the CPU's report of AVX2/FMA \\
(C3) integration & \stproved{} \stnum & the glue proposition; $31$ units in both builds, $12$ on the GPU; assumes (A2) \\
GPU filter & \stproved & Theorem E; assumes the device compiler and device correct (tested on the device at start-up), host and device code of one format (checked at start-up) \\ \midrule
compiler correctness & \stassume & tiny-prime test, arithmetic checks with two GCC versions, oracle test, start-up check; docs/BUILDING.md names the checked compilers \\
(S2), (S3) & \stassume & as above, on x86-64 with GCC~15.2 and~16.2; in every build the start-up check switches the filter off if the compiled loads or reductions disagree with exact values \\
(H) & \stassume & the enumeration's search region (the GetDecics author's statement); no overflow on $733{,}855$ real candidates and shifts, largest value $2^{59.45}$ against $2^{63}$; none among the $14{,}846{,}153{,}955$ polynomials of the $31$ test work units \\
(A2) & \stassume & no other thread of the program calls the filter \\
\bottomrule
\end{tabularx}
\end{center}

\paragraph{Open items.}
\begin{itemize}
\item Whether the size limits $8385$ and $1{,}097{,}010{,}265$ are sharp (no shipped prime is near them).  The directed search narrows the first failure to
  $(8385,15901]$ and $(1{,}097{,}010{,}265,\ 1{,}749{,}999{,}961]$ for primes $\equiv1\pmod4$; closing the gap needs a sharper proof, not more search.
\item (H) is measured on real candidates (about $12$ times headroom) and checked on every polynomial of the $31$ test work units, not proved from the search bounds.
\item The GPU filter was tested on the CPU and on one NVIDIA card; AMD, Intel and MUSA devices, and the CUDA kernel built by \code{nvcc}, were not tested.  The program
  checks the filter on every device at start-up.
\item The searches use $S=\{2,5\}$; other sets were tested on fewer polynomials.
\end{itemize}

\appendix
%======================================================================
\section{Reproducing the checks}\label{app:reproduce}

The checks are scripts of the repository, in \code{pdfilter\_Test} unless stated; they build their programs with \code{build-linux.sh} (Linux or WSL, with gcc,
gfortran and PARI).  \code{DIR} is GetDecics' \code{GetDecicsSrc} directory with the patch applied (\code{build/prepare-src.sh}).

\begin{center}\footnotesize
\begin{tabularx}{\textwidth}{@{}>{\raggedright\arraybackslash}p{0.45\textwidth}X@{}}
\toprule
command & checks \\ \midrule
\code{bash tiny\_checks.sh} (\code{QUICK=1}: a few minutes) & the compiled filter on complete sets at tiny primes \\
\code{bash arith\_checks.sh} (\code{QUICK=1}: about a minute) & the compiled reductions; the bound propagation of Theorem~\ref{thm:limits} \\
\code{bash fault\_checks.sh 100 test\_pdfilter\_oracle.cpp} & the oracle test with every deliberate fault and the moved primes \\
\code{UPSTREAM=DIR CAND=... HITS=... bash poldisc\_checks.sh} & Table~\ref{tab:poldisc}; family \code{wrap}: the constructed overflow \\
\code{TEST=test\_gpu\_filter.cpp bash build-linux.sh}, then \code{build/test\_gpu\_filter 6000000 11 8} & Theorem E at the kernel's primes \\
\code{gcc -O2 -o limit\_probe limit\_probe.c}; \code{./limit\_probe small Q SEED N climb STEPS x} & the directed search above the size limits \\
\code{python3 checks.py} & the mathematics (resultants, both chains, the size-limit computation) with independent exact models \\
\code{build/test-suite-unix.sh PROGRAM DATADIR} (repository root) & the $31$ test work units against their truth files \\
\bottomrule
\end{tabularx}
\end{center}

\paragraph{Versions measured.}  GetDecics at \code{638ca4e}.  Floating-point filter: CPU \code{327b04f}, GPU \code{718bab1}.  Montgomery CPU filter with $6+24$ primes
and scalar loads: \code{27d917f}; with vectorised loads and $24+6$ primes: \code{4094af4}.  GPU kernel with $14$ primes: \code{e758190}; with the wider table:
\code{0f874dc}; with the start-up self-test: \code{2c351c1}.

\paragraph{Not in the repository.}  The captured candidates, the Linux build recipe of the application with BOINC replaced by stand-ins, the Windows harness of the
GPU runs, the Barrett version, and the instrumented copies and scripts used for the timings and for the overflow census of the test work units.

{\small
\begin{thebibliography}{9}
\bibitem{gd} GetDecics, the NumberFields@home application, \url{https://github.com/drivere/get-decics-numberfields} (commit \code{638ca4e}).
\bibitem{cohen} H.~Cohen, \emph{A Course in Computational Algebraic Number Theory}, GTM 138, Springer, 1993.
\bibitem{lang} S.~Lang, \emph{Algebra}, revised third edition, GTM 211, Springer, 2002.
\end{thebibliography}}

\end{document}
